\section{Travaux liés}
\label{sec:related_work}

\paragraph{Réseaux convolutifs pour la classification d'images.}
Les architectures convolutives profondes, dont ResNet~\cite{he2016resnet},
ont longtemps constitué l'approche dominante en classification d'images.
Leur biais inductif de localité (les filtres convolutifs traitent des
voisinages de pixels) et d'invariance par translation les rend
particulièrement efficaces sur des jeux de données de taille modeste,
puisque ce biais réduit l'espace des hypothèses que le modèle doit
apprendre directement à partir des données.

\paragraph{Vision Transformers.}
Dosovitskiy et al.~\cite{dosovitskiy2020vit} introduisent le Vision
Transformer (ViT), qui découpe une image en une grille de patches traités
comme une séquence de tokens, à la manière d'un Transformer pour le texte.
Contrairement aux CNN, le ViT ne possède aucun biais inductif de localité
intégré à l'architecture : les relations spatiales entre patches doivent
être entièrement apprises à partir des données via les mécanismes
d'auto-attention. Les auteurs montrent que le ViT peut surpasser les CNN
à condition d'être pré-entraîné sur des jeux de données massifs
(centaines de millions d'images), mais qu'il performe nettement moins
bien que les CNN lorsqu'il est entraîné from scratch sur des jeux de
données de taille modeste — le ViT est ainsi qualifié de plus
\emph{data-hungry}.

\paragraph{DeiT et l'efficacité en données.}
Touvron et al.~\cite{touvron2021deit} (DeiT) abordent directement cette
limite en proposant une stratégie d'entraînement — notamment une
distillation par un modèle enseignant convolutif — permettant d'entraîner
des ViT compétitifs sur ImageNet sans recourir à un pré-entraînement sur
des données massives. Ce travail est particulièrement pertinent pour ce
projet : il documente empiriquement le comportement data-hungry des ViT
et propose des pistes pour l'atténuer, ce qui motive directement l'analyse
de sensibilité à la quantité de données menée en
Section~\ref{sec:ablation} (courbe d'accuracy en fonction de la fraction
du jeu d'entraînement utilisée).

\paragraph{Classification fine-grained.}
La classification fine-grained pose des difficultés spécifiques par
rapport à la classification générique : les catégories partagent une
structure globale similaire et ne se distinguent que par des détails
localisés. Le jeu de données CUB-200-2011~\cite{wah2011cub}, utilisé dans
ce projet, est une référence classique pour évaluer les architectures sur
ce type de tâche.

\paragraph{Positionnement de ce projet.}
Ce projet se situe à l'intersection de ces travaux : il reproduit à petite
échelle une comparaison ViT vs CNN sur une tâche fine-grained, en mettant
l'accent sur la sensibilité de chaque architecture à la quantité de
données disponible plutôt que sur la seule performance finale.